\documentclass[12pt,a4paper]{article}
\usepackage{cmap} % Makes the PDF copiable. See http://tex.stackexchange.com/a/64198/25761
\usepackage[T1]{fontenc}
\usepackage[brazil]{babel}
\usepackage[utf8]{inputenc}
\usepackage{amsmath}
\usepackage{amsfonts}
\usepackage{amssymb}
\usepackage{amsthm}
\usepackage{textcomp} % \degree
\usepackage{gensymb} % \degree
\usepackage[usenames,svgnames,dvipsnames]{xcolor}
\usepackage{hyperref}
\usepackage{multicol}
\usepackage{graphicx}
\usepackage[margin=2cm]{geometry}

\hypersetup{
    colorlinks = true,
    allcolors = {blue}
}

\newcommand{\fixme}{{\color{red}(...)}}

\newcommand*\sen[1]{\operatorname{sen}{#1}}
\newcommand*\abs[1]{\left|#1\right|}

\newcommand*\tipo{PROVA I}
\newcommand*\turma{ELE141-01U}
\newcommand*\disciplina{CDI0001}
\newcommand*\eu{Helder G. G. de Lima}
\newcommand*\data{16/09/2015}

\author{\eu}
\title{\tipo - \disciplina}
\date{\data}

\begin{document}
\thispagestyle{empty}
\newgeometry{margin=2cm,bottom=0.5cm}
\begin{center}
\includegraphics[width=9.0cm]{marca}
\noindent\begin{tabular}{l c c r}
  \textbf{\disciplina (\turma)}
& \textbf{\tipo}
& \textbf{\data}
\end{tabular}
\\ Prof. \eu\footnote{
Este é um material de acesso livre distribuído sob os termos da licença \href{https://creativecommons.org/licenses/by-sa/4.0/deed.pt_BR}{Creative Commons BY-SA 4.0}.}
\end{center}

\noindent Nome do(a) aluno(a): \rule{13cm}{0.01cm}

%\section*{Instruções}
\begin{center}\fbox{
\begin{minipage}{14cm}

{\footnotesize
\begin{itemize}
\renewcommand{\theenumi}{\Roman{enumi}}
\item Identifique-se em todas as folhas.
\item Mantenha o celular e os demais equipamentos eletrônicos desligados durante a prova.
\item Justifique cada resposta com cálculos ou argumentos baseados na teoria estudada.
\item Escolha \textsc{\textbf{uma}} das 6 questões para \textsc{\textbf{não}} fazer (ela não será corrigida): \framebox(30,10){}

\end{itemize}
}

\end{minipage}
}
\end{center}


%\section*{Questões}

\begin{enumerate}
\item Seja $g: D \to \mathbb{R}$ a função dada por $g(x) = x\sqrt{x^2} + \ln{\abs{x + 1}} - \ln{\abs{x - 1}}$. Explique:
\begin{enumerate}
\item (0,5) Qual é o domínio $D$ em que $g$ está definida?
\item (0,5) A função $g$ é par, ímpar ou nenhuma das duas coisas?
\item (0,5) A função $g$ é contínua nos pontos em que está definida? Por quê?
\item (0,5) Quanto é $\lim\limits_{x\to1^+} g(x)$?
\end{enumerate}

\item Seja $f(x) = \dfrac{x^2 - x - 2}{3x - 6}$ e considere $\varepsilon > 0$.
\begin{enumerate}
\item (1,5) Obtenha $\delta > 0$ (em função de $\varepsilon$) tal que ``se $0 < \abs{x - 2} < \delta$ então $\abs{f(x) - 1} < \varepsilon$''.
\item (0,5) Interprete isso em termos de limites.
\end{enumerate}

\item (2,0) Determine um intervalo $[n, n+1]$ contendo (pelo menos) uma solução positiva da equação $x^5 - 200x = 100$. O que garante que há realmente uma solução no intervalo encontrado?


\item Seja $f: \mathbb{R} \to \mathbb{R}$ a função definida por
\[
f(x) =
\left\{\begin{array}{ll}
\dfrac{\sqrt{6-x}-2}{\sqrt{3-x}-1}, & \text{se}\quad x < 2 \\
1/2, & \text{se}\quad x = 2 \\
\dfrac{\sen{(x-2)}}{\sen{(x-2)} + x-2}, & \text{se}\quad x>2.
\end{array}\right.
\]
\begin{enumerate}
\item (1,5) Calcule, se existirem, $\lim\limits_{x\to 2^-} f(x)$ e $\lim\limits_{x\to 2^+} f(x)$.
\item (0,5) Explique se $f$ é contínua em $x = 2$.
\end{enumerate}

\item Determine, se existirem:
\begin{multicols}{2}
\begin{enumerate}
\item (1,0) $\lim\limits_{x\to 1} \dfrac{\ln(x)}{x-1}$
\item (1,0) $\lim\limits_{x\to +\infty} \dfrac{\sen(x^2)}{\sqrt{x}}$
\end{enumerate}
\end{multicols}

\item Calcule, se existirem:
\begin{enumerate}
\item (1,0) $\lim\limits_{x\to 0} \dfrac{e^{3x}-1}{e^{7x}-1}$
\item (1,0) $\lim\limits_{x\to +\infty} x(\sqrt{x^2 + \pi} - x)$
\end{enumerate}
\end{enumerate}

\newpage
\restoregeometry
\section*{Respostas e observações}

\begin{enumerate}
\item \textit{Seja $g: D \to \mathbb{R}$ a função dada por $g(x) = x\sqrt{x^2} + \ln{\abs{x + 1}} - \ln{\abs{x - 1}}$. Explique:}
\begin{enumerate}
\item \textit{Qual é o domínio $D$ em que $g$ está definida?}
\begin{align*}
D
& = \operatorname{Dom}(g) \\
& = \left\{ x \in \mathbb{R} | x^2 \geq 0 \text{ e } \abs{x+1} > 0 \text{ e } \abs{x-1} > 0 \right\} \\
& = \left\{ x \in \mathbb{R} | \abs{x+1} \neq 0 \text{ e } \abs{x-1} \neq 0 \right\} \\
& = \left\{ x \in \mathbb{R} | x+1 \neq 0 \text{ e } x-1 \neq 0 \right\} \\
& = \left\{ x \in \mathbb{R} | x \neq -1 \text{ e } x\neq 1 \right\} \\
& = \mathbb{R} \setminus \{-1, 1 \} \\
& = (-\infty, -1) \cup ( -1, 1) \cup ( 1, +\infty)
\end{align*}

\item \textit{A função $g$ é par, ímpar ou nenhuma das duas coisas?}

Tem-se
\begin{align*}
g(-x)
& = (-x)\sqrt{(-x)^2} + \ln{\abs{(-x) + 1}} - \ln{\abs{(-x) - 1}} \\
& = -x \sqrt{x^2} + \ln{\abs{-(x - 1)}} - \ln{\abs{-(x+1)}} \\
& = -x \sqrt{x^2} + \ln{\abs{(x - 1)}} - \ln{\abs{(x+1)}} \\
& = -( x \sqrt{x^2} - \ln{\abs{(x - 1)}} + \ln{\abs{(x+1)}} ) \\
& = -g(x).
\end{align*}
Isso significa que $g$ é uma função ímpar.

Os cálculos também poderiam ser feitos depois de escrever que $g(x) = x\abs{x} + \ln{\frac{ \abs{x + 1} }{ \abs{x - 1} } }$.

\item \textit{A função $g$ é contínua nos pontos em que está definida? Por quê?}

Uma vez que raízes quadradas, polinômios e logaritmos são funções contínuas, e que a soma, subtração, multiplicação e composição de funções contínuas resulta em funções contínuas, pode-se concluir que a função $g$ também é contínua (pois é uma soma/composição de funções contínuas).


\item \textit{Quanto é $\lim\limits_{x\to1^+} g(x)$?}

Levando em conta que
\[
\lim\limits_{x\to1^+} x\sqrt{x^2} + \ln{\abs{x + 1}}
= 1 \sqrt{1^2} + \ln{\abs{1 + 1}}
= 1 + \ln{2} \in \mathbb{R}
\]
e que
\[
\lim\limits_{x\to1^+} - \ln{\abs{x - 1}}
= \lim\limits_{u\to 0^+} - \ln{\abs{u}}
= \lim\limits_{u\to 0^+} - \ln{u}
= +\infty,
\]
conclui-se que
\[
\lim\limits_{x\to1^+} g(x) = + \infty.
\]

\end{enumerate}

\item \textit{Seja $f(x) = \dfrac{x^2 - x - 2}{3x - 6}$ e considere $\varepsilon > 0$.}
\begin{enumerate}
\item \textit{Obtenha $\delta > 0$ (em função de $\varepsilon$) tal que ``se $0 < \abs{x - 2} < \delta$ então $\abs{f(x) - 1} < \varepsilon$''.}

Observe que
\begin{align*}
\abs{f(x) - 1} < \varepsilon
& \Leftrightarrow \abs{ \frac{x^2 - x - 2}{3x - 6} - 1} < \varepsilon
  \Leftrightarrow \abs{\frac{(x^2 - x - 2) - (3x - 6)}{3x - 6}} < \varepsilon \\
& \Leftrightarrow \abs{\frac{x^2 - 4x + 4}{3(x - 2)}} < \varepsilon
  \Leftrightarrow \abs{\frac{(x - 2)^2}{x - 2}} < 3\varepsilon
  \Leftrightarrow \abs{x - 2} < 3 \varepsilon.
\end{align*}
Portanto, para um $\varepsilon>0$ dado, basta escolher qualquer $\delta \leq 3\varepsilon$ e será verdade que ``se $0 < \abs{x - 2} < \delta$ então $\abs{f(x) - 1} < \varepsilon$''.

\item \textit{Interprete isso em termos de limites.}

O significado da propriedade anterior é que a função $f(x)$ se aproxima de $1$ quando $x$ tende a $2$, isto é, $\lim\limits_{x\to2} f(x) = 1$.

\end{enumerate}

\item \textit{Determine um intervalo $[n, n+1]$ contendo (pelo menos) uma solução positiva da equação $x^5 - 200x = 100$. O que garante que há realmente uma solução no intervalo encontrado?}

Seja $f(x) = x^5 - 200x$. Atribuindo alguns valores à variável $x$ (positivos, já que o enunciado não pede soluções negativas), observa-se que
\begin{align*}
f(1) & = 1^5 - 200 \cdot 1 = 1 - 200 = -199 < 100 \\
f(2) & = 2^5 - 200 \cdot 2 = 32 - 400 = -368 < 100 \\
\bf{ f(3) } & = 3^5 - 200 \cdot 3 = 243 - 600 =-357 \bf{ < 100 }\\
\bf{ f(4) } & = 4^5 - 200 \cdot 4 = 1024 - 800 =224 \bf{ > 100 }
\end{align*}
Como $f$ é uma função contínua (é um polinômio) e $f(3) < 100 < f(4)$, o teorema do valor intermediário garante que existe algum $c \in (3,4)$ tal que $f(c) = c^5 - 200 \cdot c = 100$, isto é, $c$ é uma solução da equação indicada.

Um raciocínio similar poderia ser feito com a função contínua $g(x) = x^5 - 200x - 100$, à qual seria aplicado o Teorema de Bolzano.

\item \textit{Seja $f: \mathbb{R} \to \mathbb{R}$ a função definida por}
\[
f(x) =
\left\{\begin{array}{ll}
\dfrac{\sqrt{6-x}-2}{\sqrt{3-x}-1}, & \text{se}\quad x < 2 \\
1/2, & \text{se}\quad x = 2 \\
\dfrac{\sen{(x-2)}}{\sen{(x-2)} + x-2}, & \text{se}\quad x > 2.
\end{array}\right.
\]
\begin{enumerate}
\item \textit{Calcule, se existirem, $\lim\limits_{x\to 2^-} f(x)$ e $\lim\limits_{x\to 2^+} f(x)$.}
Para o primeiro limite, temos:
\begin{align*}
\lim\limits_{x\to 2^-} f(x)
& = \lim\limits_{x\to 2^-} \frac{\sqrt{6-x}-2}{\sqrt{3-x}-1}
  = \lim\limits_{x\to 2^-} \frac{(\sqrt{6-x}-2)(\sqrt{6-x}+2)}{(\sqrt{3-x}-1)(\sqrt{6-x}+2)} \\
& = \lim\limits_{x\to 2^-} \frac{((6-x)-4)}{(\sqrt{3-x}-1)(\sqrt{6-x}+2)}
                           \frac{(\sqrt{3-x}+1)}{(\sqrt{3-x}+1)} \\
& = \lim\limits_{x\to 2^-} \frac{(2-x)(\sqrt{3-x}+1)}{((3-x)-1)(\sqrt{6-x}+2)}
  = \lim\limits_{x\to 2^-} \frac{(2-x)(\sqrt{3-x}+1)}{(2-x)(\sqrt{6-x}+2)} \\
& = \lim\limits_{x\to 2^-} \frac{\sqrt{3-x}+1}{\sqrt{6-x}+2}
  = \frac{\sqrt{3-2}+1}{\sqrt{6-2}+2}
  = \frac{\sqrt{1}+1}{\sqrt{4}+2}
  = \frac{1}{2}
\end{align*}

Para o segundo, temos
\begin{align*}
\lim\limits_{x\to 2^+} f(x)
& = \lim\limits_{x\to 2^+} \frac{\sen{(x-2)}}{\sen{(x-2)} + x-2}
  = \lim\limits_{u\to 0^+} \frac{\sen{u}}{\sen{u} + u}
  = \lim\limits_{u\to 0^+} \frac{\frac{\sen{u}}{u}}{\frac{\sen{u}}{u} + \frac{u}{u}}
  = \frac{1}{1 + 1}
  = \frac{1}{2}. \\
\end{align*}

\item \textit{Explique se $f$ é contínua em $x = 2$.}

Uma vez que os limites laterais $\lim\limits_{x\to 2^-} f(x)$ e $\lim\limits_{x\to 2^+} f(x)$ existem e são iguais a $1/2$, resulta que o limite bilateral $\lim\limits_{x\to 2} f(x)$ também existe e é igual a $1/2$. Além disso, como $\lim\limits_{x\to 2} f(x) = 1/2 = f(2)$, conclui-se que $f$ é uma função contínua em $x=2$.

\end{enumerate}

\item \textit{Determine, se existirem:}
\begin{enumerate}
\item
\begin{align*}
\lim\limits_{x\to 1} \frac{\ln(x)}{x-1}
& = \lim\limits_{u\to 0} \frac{\ln(1+u)}{u}
  = \lim\limits_{u\to 0} \frac{1}{u} \ln(1+u)
  = \lim\limits_{u\to 0} \ln\left((1+u)^\frac{1}{u}\right)\\
& = \ln \left( \lim\limits_{u\to 0} (1+u)^\frac{1}{u}\right)
  = \ln(e)
  = 1.
\end{align*}

Solução alternativa, considerando $v = \ln(x)$ e $x=e^v$:

\[
\lim\limits_{x\to 1} \frac{\ln(x)}{x-1}
  = \lim\limits_{v\to 0} \frac{v}{e^v-1}
  = \lim\limits_{v\to 0} \left( \frac{e^v-1}{v} \right)^{-1}
  = \left( \lim\limits_{v\to 0} \frac{e^v-1}{v} \right)^{-1}
  = (\ln(e))^{-1}
  = 1.
\]


\item $\lim\limits_{x\to +\infty} \dfrac{\sen(x^2)}{\sqrt{x}} = \lim\limits_{x\to +\infty} \dfrac{1}{\sqrt{x}} \cdot \sen(x^2) = 0$, pois $\lim\limits_{x\to +\infty} \dfrac{1}{\sqrt{x}} = 0$ e $-1 \leq \sen(x^2) \leq 1$, isto é, a função limitada $\sen(x^2)$ está sendo multiplicada por uma que tende a zero, o que implica que o produto também tende a zero.

\end{enumerate}

\item \textit{Calcule, se existirem:}
\begin{enumerate}
\item
\[
\lim\limits_{x\to 0} \frac{e^{3x}-1}{e^{7x}-1}
  = \lim\limits_{x\to 0} \frac{3x \cdot \frac{e^{3x}-1}{3x}}{7x \cdot \frac{e^{7x}-1}{7x}}
  = \frac{3}{7}\lim\limits_{x\to 0} \frac{\frac{e^{3x}-1}{3x}}{\frac{e^{7x}-1}{7x}}
  = \frac{3}{7}\frac{\lim\limits_{x\to 0} \frac{e^{3x}-1}{3x}}{\lim\limits_{x\to 0} \frac{e^{7x}-1}{7x}}
  = \frac{3}{7}\frac{\lim\limits_{u\to 0} \frac{e^{u}-1}{u}}{\lim\limits_{v\to 0} \frac{e^{v}-1}{v}}
  = \frac{3}{7} \cdot \frac{1}{1}
  = \frac{3}{7}\\
\]

Solução alternativa:
\[
\lim\limits_{x\to 0} \frac{e^{3x}-1}{e^{7x}-1}
  = \lim\limits_{x\to 0} \frac{(e^3)^x-1}{(e^7)^x-1}
  = \lim\limits_{x\to 0} \frac{\frac{(e^3)^x-1}{x}}{\frac{(e^7)^x-1}{x}}
  = \frac{\ln(e^3)}{\ln(e^7)}
  = \frac{3 \ln(e)}{7\ln(e)}
  = \frac{3}{7}.
\]

Outra solução, considerando $u=e^x$:
\begin{align*}
\lim\limits_{x\to 0} \frac{e^{3x}-1}{e^{7x}-1}
& = \lim\limits_{x\to 0} \frac{(e^x)^3-1}{(e^x)^7-1}
  = \lim\limits_{u\to 1} \frac{u^3-1}{u^7-1} \\
& = \lim\limits_{u\to 1} \frac{(u-1)(u^2 + u + 1)}{(u-1)(u^6 + u^5 + u^4 + u^3 + u^2 + u + 1)} \\
& = \lim\limits_{u\to 1} \frac{(u^2 + u + 1)}{(u^6 + u^5 + u^4 + u^3 + u^2 + u + 1)} \\
& = \frac{(1^2 + 1 + 1)}{(1^6 + 1^5 + 1^4 + 1^3 + 1^2 + 1 + 1)}
  = \frac{3}{7}.
\end{align*}

\item
\begin{align*}
\lim\limits_{x\to +\infty} x(\sqrt{x^2 + \pi} - x)
& = \lim\limits_{x\to +\infty} x(\sqrt{x^2 + \pi} - x) \cdot \frac{\sqrt{x^2 + \pi} + x}{\sqrt{x^2 + \pi} + x} \\
& = \lim\limits_{x\to +\infty} \frac{x((x^2 + \pi) - x^2)}{\sqrt{x^2 + \pi} + x}
  = \lim\limits_{x\to +\infty} \frac{\pi x}{|x|\sqrt{1 + \pi/x^2} + x} \\
& = \lim\limits_{x\to +\infty} \frac{\pi x}{x (\sqrt{1 + \pi/x^2} + 1)}
  = \frac{\pi}{\sqrt{1 + 0} + 1}
  = \frac{\pi}{2}
\end{align*}
\end{enumerate}
\end{enumerate}
\end{document}
